\documentclass[11pt, twocolumn, landscape, a4paper]{article}

\usepackage[margin=0.9cm]{geometry}
\usepackage[utf8]{inputenc}
\usepackage{amssymb, amsmath, amsbsy}
\usepackage{mathpazo}
\usepackage{verbatim}
\usepackage{enumerate}
\usepackage[pdftex]{hyperref}
\hypersetup{pdftitle={Vectores, Producto punto, Norma de un vector},pdfauthor={Luis Carrillo Gutiérrez}}

\renewcommand{\familydefault}{\sfdefault}

\begin{document}

\subsection*{Multiplicación de vector por un escalar}
\noindent Sea el vector $\vec{A} = (a_1,\,a_2) \in \mathbb{R}^2$ 
y sea $s \in \mathbb{R}$, se tiene : \newline 
\begin{eqnarray*}
s \cdot \vec{A} &=& s \cdot (a_1,\;a_2) \\
s\vec{A} &=&  (s \cdot a_1,\; s \cdot a_2) \in \mathbb{R}^2
\end{eqnarray*}

% Producto punto
\subsection*{Producto interior o producto escalar de un vector}
Dados los vectores $\vec{A} = (a_1,\;a_2)$ y $\vec{B} = (b_1,\;b_2) \in \mathbb{R}^2$, el producto interior o escalar de $\vec{A}$ con $\vec{B}$, es el número real $\vec{A}\cdot\vec{B}$, definido por:  

\begin{eqnarray*}
\vec{A} \cdot \vec{B} &=& (a_{1},\;a_{2}) \cdot (b_{1},\;b_{2}) \\ 
\empty &=& a_{1}b_{1} + a_{2}b_{2}\in \mathbb{R}
\end{eqnarray*}

\paragraph*{Propiedades} : \quad$\forall\;\vec{A},\,\vec{B} \in \mathbb{R}^{2}$ y $k \in \mathbb{R}$
\begin{enumerate}[i.]
\item{$\vec{A} \cdot \vec{B} \in \mathbb{R}$}
\item{$\vec{A} \cdot \vec{B} = \vec{B} \cdot \vec{A}$}
\item{$(k\,\vec{A}) \cdot \vec{B} = \vec{A} \cdot (k\,\vec{B}) = k\,(\vec{A} \cdot \vec{B})$}
\item{$\vec{A} \cdot (\vec{B} + \vec{C}) = \vec{A} \cdot \vec{B} + \vec{A} \cdot \vec{C}$}
\item{$\vec{A} \cdot \vec{A} \ge 0$, si $\vec{A} = (a_{1},\,a_{2}) \implies \vec{A}\cdot \vec{A} = {a_{1}}^{2} + {a_{2}}^{2} > 0$}
\end{enumerate}

\subsection*{Norma o longitud de un vector}
Sea $\vec{A} \in \mathbb{R}$, la norma o longitud del vector $\vec{A}$, es un número real positivo $[|\,\vec{A}\,|]$, definido por :  
\[
\Big{||} \vec{A} || = \sqrt{\vec{A} \cdot \vec{A}\,}, \quad\text{si}\quad\vec{A} = (a_{1},\,a_{2}) \implies || \vec{A} || = \sqrt{{a_{1}}^{2} + {a_{2}}^{2}\,}
\]
\paragraph*{Propieades} : \quad$\forall\;\vec{A} \in \mathbb{R}^{2}$ y $k \in \mathbb{R}$ 
%Ak AA
% \subsubsection*{Propiedades}
\begin{enumerate}[i.]
\item{$||\vec{A}|| \in \mathbb{R}_{0}^{+}$
\begin{itemize}
\item{$||\vec{A}|| \ge 0$}
\item{$||\vec{A}|| = 0 \iff \vec{A} = \vec{0}$}
\end{itemize}
}
\item{$||\vec{A}||^{\,2} = \vec{A} \cdot \vec{A}$}
\item{$||\vec{A}|| = ||- \vec{A}||$}
\item{$||k\,\vec{A}|| = |k|\,||\,\vec{A}\,||$}
\item{$||\vec{A} + \vec{B}\,|| \le ||\vec{A}|| + ||\vec{B}||$\qquad Desigualdad triangular}
\item{$||\vec{A} \cdot \vec{B}\,|| \le ||\vec{A}|| \cdot ||\vec{B}||$\qquad Desigualdad de Cauchy-Schwarz}
\end{enumerate}

\end{document}